% Options for packages loaded elsewhere
\PassOptionsToPackage{unicode}{hyperref}
\PassOptionsToPackage{hyphens}{url}
\documentclass[
  ignorenonframetext,
]{beamer}
\newif\ifbibliography
\usepackage{pgfpages}
\setbeamertemplate{caption}[numbered]
\setbeamertemplate{caption label separator}{: }
\setbeamercolor{caption name}{fg=normal text.fg}
\beamertemplatenavigationsymbolsempty
% remove section numbering
\setbeamertemplate{part page}{
  \centering
  \begin{beamercolorbox}[sep=16pt,center]{part title}
    \usebeamerfont{part title}\insertpart\par
  \end{beamercolorbox}
}
\setbeamertemplate{section page}{
  \centering
  \begin{beamercolorbox}[sep=12pt,center]{section title}
    \usebeamerfont{section title}\insertsection\par
  \end{beamercolorbox}
}
\setbeamertemplate{subsection page}{
  \centering
  \begin{beamercolorbox}[sep=8pt,center]{subsection title}
    \usebeamerfont{subsection title}\insertsubsection\par
  \end{beamercolorbox}
}
% Prevent slide breaks in the middle of a paragraph
\widowpenalties 1 10000
\raggedbottom
\AtBeginPart{
  \frame{\partpage}
}
\AtBeginSection{
  \ifbibliography
  \else
    \frame{\sectionpage}
  \fi
}
\AtBeginSubsection{
  \frame{\subsectionpage}
}
\usepackage{iftex}
\ifPDFTeX
  \usepackage[T1]{fontenc}
  \usepackage[utf8]{inputenc}
  \usepackage{textcomp} % provide euro and other symbols
\else % if luatex or xetex
  \usepackage{unicode-math} % this also loads fontspec
  \defaultfontfeatures{Scale=MatchLowercase}
  \defaultfontfeatures[\rmfamily]{Ligatures=TeX,Scale=1}
\fi
\usepackage{lmodern}
\ifPDFTeX\else
  % xetex/luatex font selection
\fi
% Use upquote if available, for straight quotes in verbatim environments
\IfFileExists{upquote.sty}{\usepackage{upquote}}{}
\IfFileExists{microtype.sty}{% use microtype if available
  \usepackage[]{microtype}
  \UseMicrotypeSet[protrusion]{basicmath} % disable protrusion for tt fonts
}{}
\makeatletter
\@ifundefined{KOMAClassName}{% if non-KOMA class
  \IfFileExists{parskip.sty}{%
    \usepackage{parskip}
  }{% else
    \setlength{\parindent}{0pt}
    \setlength{\parskip}{6pt plus 2pt minus 1pt}}
}{% if KOMA class
  \KOMAoptions{parskip=half}}
\makeatother
\usepackage{graphicx}
\makeatletter
\newsavebox\pandoc@box
\newcommand*\pandocbounded[1]{% scales image to fit in text height/width
  \sbox\pandoc@box{#1}%
  \Gscale@div\@tempa{\textheight}{\dimexpr\ht\pandoc@box+\dp\pandoc@box\relax}%
  \Gscale@div\@tempb{\linewidth}{\wd\pandoc@box}%
  \ifdim\@tempb\p@<\@tempa\p@\let\@tempa\@tempb\fi% select the smaller of both
  \ifdim\@tempa\p@<\p@\scalebox{\@tempa}{\usebox\pandoc@box}%
  \else\usebox{\pandoc@box}%
  \fi%
}
% Set default figure placement to htbp
\def\fps@figure{htbp}
\makeatother
\setlength{\emergencystretch}{3em} % prevent overfull lines
\providecommand{\tightlist}{%
  \setlength{\itemsep}{0pt}\setlength{\parskip}{0pt}}
\usepackage{bookmark}
\IfFileExists{xurl.sty}{\usepackage{xurl}}{} % add URL line breaks if available
\urlstyle{same}
\hypersetup{
  hidelinks,
  pdfcreator={LaTeX via pandoc}}

\author{\texorpdfstring{}{}}
\date{}

\begin{document}

\section{Introduction: Case Diagrams and Categorical
Linguistics}\label{sec:introduction}

\begin{frame}{The Architecture of Case Assignment}
\protect\phantomsection\label{the-architecture-of-case-assignment}
Natural language encodes situational relations---\emph{who does what to
whom, with what, where, and why}---through various mechanisms, including
noun case systems. In morphologically rich languages like Finnish,
Latin, or Sanskrit, this case encoding is overt: distinct affixes mark
the agent, patient, instrument, and location. In languages like Mandarin
or English, this relational skeleton is orchestrated through word order
and adpositions. Beyond the surface variation of morphological case, the
underlying computational problem is universal: a cognitive agent
benefits from assembling and assigning \emph{structured roles} to the
participants of every scenario it encounters ()), predicting and
sketching relationships against the constraints of a conceptual plane.

This universality suggests that case is not merely a morphosyntactic
accident, but a reflection of a deeper, embodied cognitive architecture.
This paper reviews the long lineage of this hypothesis, tracing its
roots from Pāṇini's ancient \emph{kāraka} theory (which first abstracted
nominal roles derived from actions), through Roman Jakobson's
structuralist feature matrices that viewed cases as oppositional
spatial/semantic networks, to Fillmore's {[}-@fillmore1968case{]} ``Case
for Case'' which posited a universal inventory of \emph{deep cases}
(Agent, Patient, Instrument). We review how Igor Mel'čuk's Meaning-Text
Theory integrated these insights, treating cases as profound relational
networks mapping form to deep semantic substance. Subsequent typological
work by Polinsky and Preminger {[}-@polinsky2015case{]}, Blake
{[}-@blake2001grammatical{]}, and Haspelmath
{[}-@haspelmath2009universality{]} has confirmed that the space of
variation across human language is tightly structured into a small
number of \emph{alignment types} related by systematic transformations.

This review highlights how \textbf{category theory} provides the natural
mathematical language for formalizing this structure, and how
\textbf{commutative diagrams} function not merely as convenient
notation, but as \emph{cognitively privileged, embodied
representations}. By constraining relationships to a two-dimensional
graph, these diagrams mimic the hand-eye coordination of sketching
thought, providing inferential advantages unavailable to purely linear,
sentential encodings.
\end{frame}

\begin{frame}{Why Diagrams? Constraints on the Plane}
\protect\phantomsection\label{why-diagrams-constraints-on-the-plane}
The cognitive science of diagrammatic reasoning provides a strong
empirical foundation for this claim. To cast a problem onto a
two-dimensional visual topology is to recruit the nervous system's
innate, ancient capacity for spatial navigation. Larkin and Simon
{[}-@larkin1987diagram{]} demonstrated that diagrams can be
computationally superior to sentential representations: their planar
constraint enables efficient perceptual grouping that would otherwise
require painstaking, explicit logical deduction. Shimojima
{[}-@shimojima1996reasoning{]} refined this with the concept of
\emph{free ride inferences}---conclusions that ``fall out''
automatically from the geometric constraints of a sketch without
syntactic manipulation.

In the context of case theory, commutative diagrams offer exactly these
embodied advantages. A commutative triangle expressing that morphism
\(g \circ f\) equals morphism \(h\) simultaneously encodes:

\begin{enumerate}
\tightlist
\item
  The compositional structure of the relation (two steps factor through
  an intermediate case role).
\item
  The \emph{commutativity constraint} (the direct and indirect visual
  routes yield the same mathematical result).
\item
  The full relational context (the entire shape of the scenario is
  apprehensible in a single perceptual glance).
\end{enumerate}

This is exactly the kind of gestalt understanding that linear notation
obscures. \autoref{fig:case-minimal} illustrates this with a minimal
three-case category. As Giardino {[}-@giardino2017diagrammatic{]}
emphasizes, mathematical diagrams engage a hybrid mode of reasoning,
intimately wedding perceptual pattern-recognition with theoretical
knowledge---a mode effortlessly aligned with the predictive processing
architecture of active inference {[}@friston2017active{]}. The
intellectual lineage traces through Peirce's \emph{existential graphs}:
his spatial logic system proved that first-order reasoning can be
conducted entirely diagrammatically on a ``sheet of assertion.'' Our
case diagrams extend this Peircean tradition into relational semantics:
inference proceeds not by algebraic churning, but by tracing lines and
composing arrows---drawing the contours of thought itself.

\begin{figure}
\centering
\pandocbounded{\includegraphics[keepaspectratio,alt={A minimal case category \textbackslash mathcal\{C\} with three objects---Agent (NOM), Patient (ACC), and Instrument (INS)---and three morphisms: f\textbackslash colon\textbackslash text\{NOM\}\textbackslash to\textbackslash text\{ACC\} (transitive action, weight w=0.9), g\textbackslash colon\textbackslash text\{INS\}\textbackslash to\textbackslash text\{ACC\} (instrumental application, w=0.7), and the composite g \textbackslash circ h\textbackslash colon\textbackslash text\{NOM\}\textbackslash to\textbackslash text\{INS\}\textbackslash to\textbackslash text\{ACC\} (acts-via, w=0.9 \textbackslash times 0.7 = 0.63). The commutative triangle encodes that the direct path f and the factored path g \textbackslash circ h yield the same patient---making the argument-structure factorization through the instrument role visually immediate. This exemplifies Shimojima's {[}-@shimojima1996reasoning{]} free-ride inference: the two-dimensional layout reveals compositional structure that linear notation obscures.}]{output/figures/case_category_minimal.png}}
\caption{A minimal case category \(\mathcal{C}\) with three
objects---Agent (NOM), Patient (ACC), and Instrument (INS)---and three
morphisms: \(f\colon\text{NOM}\to\text{ACC}\) (transitive action, weight
\(w=0.9\)), \(g\colon\text{INS}\to\text{ACC}\) (instrumental
application, \(w=0.7\)), and the composite
\(g \circ h\colon\text{NOM}\to\text{INS}\to\text{ACC}\) (acts-via,
\(w=0.9 \times 0.7 = 0.63\)). The commutative triangle encodes that the
direct path \(f\) and the factored path \(g \circ h\) yield the same
patient---making the argument-structure factorization through the
instrument role visually immediate. This exemplifies Shimojima's
{[}-@shimojima1996reasoning{]} free-ride inference: the two-dimensional
layout reveals compositional structure that linear notation
obscures.}\label{fig:case-minimal}
\end{figure}
\end{frame}

\begin{frame}{Five Converging Research Pillars}
\protect\phantomsection\label{five-converging-research-pillars}
The present synthesis draws on five research traditions, each
contributing an essential formal ingredient:

\begin{block}{Pillar 1: Case Systems, Alignment Typology, and
Grammatical Relations}
\protect\phantomsection\label{pillar-1-case-systems-alignment-typology-and-grammatical-relations}
The empirical foundation comes from linguistic typology. We formalize
the case inventory and alignment systems catalogued by Polinsky and
Preminger {[}-@polinsky2015case{]}, Claassen
{[}-@claassen2025alignment{]}, and Wu {[}-@wu2024amis{]} as categories
with case roles as objects and grammatical relations as morphisms.
Dowty's {[}-@dowty1991thematic{]} proto-role theory---which decomposes
thematic roles into clusters of entailments rather than discrete
atoms---motivates our use of enriched (weighted) morphisms. \textbf{Our
novel addition, Synthetic Case-Role Algebra, extends this by formalizing
these proto-roles as objects in a {[}0,1{]}-enriched monoidal category,
enabling purely algebraic manipulation of semantic roles contextually.}
\end{block}

\begin{block}{Pillar 2: Categorial and Type-Theoretic Grammar}
\protect\phantomsection\label{pillar-2-categorial-and-type-theoretic-grammar}
Lambek's {[}-@lambek1958mathematics{]} syntactic calculus reformulates
grammatical combination as algebraic cancellation in a pregroup,
yielding derivations that can be visualized as Joyal and Street's
{[}-@joyalstreet1991geometry{]} string diagrams. Song
{[}-@song2022act{]} extends this with monadic semantics for root syntax,
showing that category-theoretic structure reaches down to the sublexical
level.
\end{block}

\begin{block}{Pillar 3: Categorical Compositional Distributional
Semantics (DisCoCat)}
\protect\phantomsection\label{pillar-3-categorical-compositional-distributional-semantics-discocat}
Coecke, Sadrzadeh, and Clark's {[}-@coecke2010mathematical{]} DisCoCat
framework composes distributional word meanings according to syntactic
structure using monoidal categories and string diagrams. This framework
bridges the two great traditions of linguistic meaning---formal
(truth-conditional, compositional) and distributional
(context-dependent, statistical)---by using the algebra of compact
closed categories to compose vector representations according to
type-logical derivations. The resulting categorical semantics provides
the \emph{algebraic formalization} of the distributional programme that
modern large language models (from Word2Vec {[}@mikolov2013efficient{]}
through transformer architectures {[}@vaswani2017attention;
@devlin2019bert{]}) implement empirically, making it possible to analyse
both symbolic and neural approaches to language within a single
mathematical framework. Recent extensions include DisCoCirc
{[}@defelice2020discourse; @defelice2022discocirc{]} for discourse-level
coherence and the lambeq library {[}@lorenz2023lambeq{]} for quantum
NLP, which implements the full categorical pipeline on quantum hardware.
\end{block}

\begin{block}{Pillar 4: Enriched Category Theory and Distributional
Measures}
\protect\phantomsection\label{pillar-4-enriched-category-theory-and-distributional-measures}
Bradley et al.'s {[}-@fritz2021enriched{]} enrichment over \([0,1]\)
equips hom-sets with distributional proximity measures. Bradley's
{[}-@bradley2020entropy; -@bradley2024ipam{]} work on topological
operads and categorical magnitude provides information-theoretic
invariants that quantify the ``effective size'' of a linguistic
category---a measure that captures how much distributional information
the category encodes. Leinster and Shulman {[}-@leinster2021magnitude{]}
extend this to \textbf{magnitude homology}, categorifying magnitude into
a graded homological invariant that detects higher-dimensional
structural differences between case systems indistinguishable by scalar
magnitude alone.
\end{block}

\begin{block}{Pillar 5: Topos-Theoretic Bridges and Inter-Theoretic
Transfer}
\protect\phantomsection\label{pillar-5-topos-theoretic-bridges-and-inter-theoretic-transfer}
Caramello's {[}-@caramello2016bridges; -@caramello2021five;
-@caramello2023syntactic{]} bridge technique uses classifying toposes as
transfer points between mathematical theories, enabling results proved
in one formalization of case to port automatically to another. Phillips
{[}-@phillips2024lot{]} strengthens this by showing that the Language of
Thought admits universal constructions in a topos, grounding the
systematicity of case assignment in the deepest structures of
categorical logic.
\end{block}
\end{frame}

\begin{frame}{Unifying Integration: Active Inference and the CEREBRUM
Architecture}
\protect\phantomsection\label{unifying-integration-active-inference-and-the-cerebrum-architecture}
The unifying framework is \textbf{active inference}
{[}@friston2017active{]}, which casts cognition as approximate Bayesian
inference under a generative model. Case-marked commutative diagrams
serve as the structural core of such generative models: each diagram
specifies a pattern of expected relational dependencies, and the agent's
task in understanding (or producing) a sentence is to minimize surprise
relative to this diagrammatic prior. The \textbf{CEREBRUM} architecture
{[}@friedman2024cerebrum{]}---Case-Enabled Reasoning Engine with
Bayesian Representations for Unified Modeling---provides a computational
instantiation of this framework, with case roles as functional
specializations of model components in an active inference cycle.

This integration is not merely metaphorical. The situation semantics of
Barwise and Perry {[}-@barwise1983situations{]} already conceptualized
meaning as structured \emph{situations} with typed constituents---an
ontology that maps naturally onto the objects and morphisms of our case
categories. Active inference adds dynamics: the agent actively samples
evidence to confirm or update its case assignments, using diagrammatic
structure to guide exploration.
\end{frame}

\begin{frame}{Paper Outline and Chapter Overview}
\protect\phantomsection\label{paper-outline-and-chapter-overview}
The remainder of the paper is organized as follows.
\autoref{sec:case-systems} reviews case systems and alignment typology,
formalizing them as categories with functors between alignment
types---including Fluid-S as a context-dependent functor parameterized
by agentive volition. \autoref{sec:categorial-grammar} develops the
categorial grammar foundations, including pregroup types, string
diagrams, and passivization as algebraic type permutation.
\autoref{sec:categorical-semantics} presents the DisCoCat framework and
its case-enriched extension, demonstrating that categorical semantics
provides the algebraic formalization of the distributional programme
underlying modern LLMs. \autoref{sec:compact-closure-discourse} extends
the framework through the compact closure axiom, diagrammatic complexity
metrics, DisCoCirc discourse circuits with case role reversal, and the
lambeq Gen II quantum NLP pipeline. \autoref{sec:enriched-categories}
introduces \([0,1]\)-enriched categories with categorical magnitude and
magnitude homology as complexity and topological invariants.
\autoref{sec:topos-theory} develops Caramello's topos-theoretic bridges
for inter-theoretic transfer via Morita equivalence, with a four-phase
algorithmic sketch for topos-theoretic grammatical induction.
\autoref{sec:cognitive-integration} integrates the five pillars within
active inference, CEREBRUM, and Distributional Active Inference,
deriving falsifiable electrophysiological predictions linking enriched
weights to ERP amplitudes. \autoref{sec:quantum-active-inference}
extends the framework into quantum topological computation via TQNNs,
the ZX-calculus, and sheaf-theoretic semantic communication.
\autoref{sec:ai-implications} develops implications for categorical deep
learning, multi-agent protocols, and compositional game theory.
\autoref{sec:cognitive-security} introduces case-theoretic firewalls and
epistemic case security as defenses against adversarial case-frame
manipulation, with a formal characterization of prompt injection as
decidable categorical type violation. \autoref{sec:conclusion} concludes
with eleven contributions and nine future research directions. The
following sections develop this roadmap in detail.
\end{frame}

\end{document}
